\section{Implementation}
\label{sec:implementation}

The MI engine is implemented in Rust ($\sim$10,000 lines, 180+ unit tests). The choice of Rust is motivated by three requirements: (i) memory safety without GC pauses, (ii) zero-cost abstractions for the inner simulation loop, and (iii) safe concurrency for parallel morphon updates. The crate compiles to native, WebAssembly (for an interactive in-browser visualizer), and Python bindings via PyO3 (for benchmark harnessing).

\subsection{Core Data Structures}

Morphons are stored in a \texttt{HashMap} keyed by stable IDs. The topology is a directed graph using the \texttt{petgraph} crate, with synapses as edge weights. We use generational indices so that references survive birth/death events.

\subsection{The Main Loop}

\begin{algorithm}
\caption{System::step (one fast tick)}
\label{alg:step}
\begin{algorithmic}[1]
\State \textbf{1.} Spike propagation: deliver pending spikes from previous step
\State \textbf{2.} Inhibitory competition (k-WTA or local interneurons)
\State \textbf{3.} Morphon update: integrate inputs, decide firing, pay metabolic cost
\State \textbf{4.} Generate new spikes from fired morphons
\If{tick.medium}
  \State \textbf{5.} Update eligibility traces, apply three-factor learning
  \State \textbf{6.} Run Endoquilibrium tick (sense vitals, adjust channels)
\EndIf
\If{tick.slow}
  \State \textbf{7.} Synaptogenesis, pruning, migration
\EndIf
\If{tick.glacial}
  \State \textbf{8.} Division, differentiation, fusion, apoptosis
\EndIf
\If{tick.homeostasis}
  \State \textbf{9.} Synaptic scaling, inter-cluster inhibition
\EndIf
\end{algorithmic}
\end{algorithm}

\subsection{Output Pathways}

MI provides two output pathways with very different characteristics. Both are biologically motivated.

\paragraph{Spike-based readout.}
The default pathway reads motor morphon potentials after spike propagation. This is the path that drives action selection in CartPole. It is the biological default --- everything happens via spikes.

\paragraph{Analog readout (Purkinje-style bypass).}
For classification tasks, MI offers an analog readout that reads associative morphon potentials \emph{directly}, bypassing spike conversion: $V_j = b_j + \sum_i w_{ji} \cdot (\sigma(p_i) - 0.5)$ where $p_i$ is the potential of associative morphon $i$, $w_{ji}$ is a learnable readout weight, and $b_j$ is a learnable per-output bias. This is biologically motivated by the Purkinje cell, which performs analog dendritic integration of climbing fiber and parallel fiber inputs to produce motor output, bypassing the spike-based intermediate computation in the cerebellar cortex.

The exact form of the readout matters. Three modifications were each individually necessary to make the readout learn:
\begin{enumerate}
  \item \textbf{Centered sigmoid} ($\sigma(p) - 0.5$ instead of $\sigma(p)$). The resting potential is zero, so $\sigma(0) = 0.5$. With $\sim$80 associative morphons, the sum of resting sigmoids is $\sim$40, which dwarfs the small state-dependent signal ($\sim$0.005 per morphon). Centering removes this constant offset and lets the weights focus on state discrimination.
  \item \textbf{Learnable per-output bias} $b_j$. Absorbs whatever mean output level the system needs, freeing the weights to encode only state differences. Without this the weights have to do double duty (mean + discrimination) and converge slowly.
  \item \textbf{No L2 weight decay}. Standard delta rule training adds $-\lambda w$ to each step. With $\lambda = 0.001$ and 200 steps per CartPole episode, decay erases $\sim$18\% of the learned weight per episode --- the gradient is consumed fighting the decay. We use weight clipping to $[-5, 5]$ for stability instead.
\end{enumerate}
Each fix individually moved CartPole from $\sim$10 steps (random baseline) toward solving; together they produced the SOLVED result. The diagnostic that surfaces this is the \emph{output discrimination} metric $d_\pm$: the difference $V_1 - V_0$ measured separately for left-leaning and right-leaning pole probes. When $d_+$ and $d_-$ have opposite signs and $|d| > 5$, the readout has learned a usable policy. With any of the three fixes missing, $d_+$ and $d_-$ remain same-signed and the policy is stuck.

The two pathways serve different purposes. Spike-based readout is closed-loop and online (suitable for control). Analog readout is open-loop and trained by a delta rule with DFA backprojection, providing exact gradients on the readout weights without requiring backpropagation through the spike network.

\subsection{Three-Factor Learning Rule}

The three-factor STDP rule on the medium path:
\begin{equation}
\Delta w_{ij} = \eta_{ij} \cdot e_{ij}(t) \cdot M_j(t) \cdot g_a
\end{equation}
where $e_{ij}$ is the eligibility trace, $M_j$ is the receptor-gated modulator signal at the post-synaptic morphon, $\eta_{ij}$ is a per-synapse plasticity rate, and $g_a$ is an astrocytic gate.

The eligibility trace integrates STDP coincidences:
\begin{equation}
\frac{de_{ij}}{dt} = -\frac{e_{ij}}{\tau_e} + A_+ \cdot \text{tr}_i^{\text{pre}} \cdot \delta_j^{\text{post}} - A_- \cdot \text{tr}_j^{\text{post}} \cdot \delta_i^{\text{pre}}
\end{equation}
with weight-dependent bounds \cite{gilson2011stability} that make weak synapses easy to strengthen and strong synapses easy to weaken.

\subsection{Tag and Capture}

We implement the synaptic tagging-and-capture model of Frey and Morris \cite{frey1997synaptic} with one critical departure: capture is \emph{episode-gated}, not per-tick. The naive per-tick rule (``capture when reward $>$ threshold'') fails on RL tasks because the reward signal saturates at 1.0 during \emph{any} successful CartPole step --- every alive step captures everything, the weight distribution collapses to uniformity (weight std drops from 1.5 to 0.15), and learning stops. The fix is to defer the capture decision to episode end and gate it on relative performance: if this episode beat the running average, capture the tagged synapses proportional to how much it beat the average; if it underperformed, decay the tags instead.

Concretely: synapses with high recent eligibility set a tag; tags accumulate at a fixed local rate (\texttt{tag\_accumulation\_rate}=0.3) and decay slowly ($\tau_{\text{tag}}=200$ steps). The \texttt{consolidation\_level} field on each synapse is a continuous $[0, 1]$ value (replacing the binary ``consolidated'' flag), and at episode end, tagged synapses gain consolidation proportional to performance delta. Consolidated synapses retain residual plasticity (10\% at \texttt{consolidation\_level}=1.0) so the network can still adapt without destroying learned features.

Endoquilibrium gates the global capture rate via a dedicated \texttt{consolidation\_gain} channel. This separation matches the biology: \emph{tags are local and continuous} (driven by per-synapse STDP eligibility), \emph{capture is global and episodic} (driven by neuromodulatory protein-synthesis-related signals during salient events). The earlier design that combined them under a single ``endo\_tag\_factor'' was biologically incorrect and produced an unstable feedback loop.

\subsection{Reproducibility}

The full benchmark suite is reproducible from a clean checkout:
\begin{verbatim}
cargo run --example cartpole --release
cargo run --example mnist_v2 --release
cargo run --example nlp_readiness --release
\end{verbatim}
Each example saves a JSON record of its configuration and results to \texttt{docs/benchmark\_results/v\{version\}/}. All numbers reported in this paper are from these JSON records.
